% chapters/abstract.tex - 中英文摘要

% ========== 中文摘要 ==========
\pagestyle{empty}
\begin{center}
\textbf{摘\quad 要}
\end{center}
\addcontentsline{toc}{chapter}{摘要}

\textbf{目的：}在规范化储备重复次数（RIR）自我指导已具备科学基础与可及性的前提下，基于计算机视觉的移动端数智化监控方案是否具备值得投入额外成本的适应性增量与实施可行性，仍缺乏针对性实证证据。本研究通过工具效度验证、预试验性随机对照训练比较与半结构化访谈三阶段递进设计，系统评估自研移动端计算机视觉系统 SportSci Pro 在青年男性下肢抗阻训练数智化监控中的测量学基础、干预方案可行性、探索性干预效果与用户接受度，并为后续正式随机对照试验提供效应量估计与流程优化参数。

\textbf{方法：}研究一招募 13 名有训练经验的健康青年男性（与研究二样本相互独立），在受控递增负荷实验条件下同步采集 SportSci Pro 与 GymAware 线性位置传感器的杠铃后蹲平均向心速度数据（60 对配对），采用组内相关系数（ICC）、Bland-Altman 分析与平均绝对误差（MAE）评估测量一致性。研究二采用两组平行、开放标签、分层区组随机对照预试验设计，将 36 名完成随机化的健康青年男性训练者（AI 辅助组 18 人，自我指导组 18 人）纳入 4 周、每周 2 次的下肢抗阻训练干预；AI 辅助组在深蹲主项中接受基于速度损失阈值的自动停组、基于纵跳准备度评估的动态减组及速度反馈负荷调控，自我指导组依据书面训练计划和储备重复次数（RIR）进行自主调节；主要考察招募保留、训练完成率、主观疲劳、Hooper 主观恢复状态、AI 辅助组速度目标命中率、自动减组执行忠实度、系统可用性量表（SUS）与技术接受度等可行性指标，并探索两组在深蹲 1RM、反向运动跳（CMJ）、蹲跳（SJ）与训练自我效能上的差异；同时同步采集 SportSci Pro 与 GymAware 的真实训练配对数据用于生态效度补充验证。T1 后测完成后对 12 名受试者（AI 组 5 人，自我指导组 7 人）进行半结构化深度访谈，采用反思性主题分析法揭示训练者的主观体验与决策机制。

\textbf{结果：}\textbf{工具效度方面}，研究一实验室条件下 SportSci Pro 与 GymAware 在 MCV 测量上达到优秀一致性（ICC = 0.940，95\% CI：0.854 至 0.975；MAE = 0.047 m/s；Bias = +0.002 m/s）。研究二真实训练场景 Rep 级配对（n = 43，3 名受试者）ICC = 0.991（Bias = +0.0019 m/s），热身级配对（n = 88，AI 辅助组全体 11 名受试者）ICC = 0.987，但存在 +0.0212 m/s 的稳定系统性正向偏移（88 对配对 100\% 均为正偏）。\textbf{干预可行性方面}，36 名受试者完成随机化，最终 24 人（AI 组 11 人，自我指导组 13 人）完成主要前后测并进入符合方案集分析。PP 样本训练完成率达 98.4\%（189/192），高依从率达 100\%（24/24），AI 组速度目标命中率达 80.2\% $\pm$ 6.9\%，自动减组触发忠实度达 100\%，严重不良事件为 0 例；9 项预设推进标准中 7 项达标，但主要后测完成率（66.7\%，24/36）与 AI 组 SUS 评分（65.00 $\pm$ 4.61 分）两项未达 $\geq$80\% 与 $\geq$70 分的预设阈值。\textbf{训练执行与内部负荷方面}，两组在全期总负荷（AI 辅助组 17651.8 $\pm$ 2198.6 kg vs. 自我指导组 17461.5 $\pm$ 2792.1 kg，$p$ = 0.854）、深蹲总负荷（$p$ = 0.988）、全期平均 sRPE（6.5 vs. 6.4）及个体 Hooper 均值（20.4 vs. 19.8）上高度对等；AI 辅助组总重复次数显著少于自我指导组（100.3 vs. 105.5 次，$p$ = 0.035，Hedges' $g$ = $-$0.95），派生的每次重复平均机械负荷高出约 6.1\%。Hooper 线性混合模型（Kenward-Roger）显示组别主效应 $p$ = 0.421、交互效应 $p$ = 0.945、课次主效应 $p$ = 0.004，两组均呈现先升后降的时间演变形态。\textbf{探索性适应结局方面}，AI 组绝对 1RM（调整后差值 +2.43 kg，95\% CI：$-$2.01 至 6.87；$g$ = 0.46）、相对 1RM（+0.03 kg/kg；$g$ = 0.39）与 SJ 高度（+0.39 cm；$g$ = 0.43）点估计均高于自我指导组但差值 95\% CI 均跨越零；CMJ 高度调整后差值 +2.16 cm（95\% CI：$-$0.01 至 4.33，$p$ = 0.051，$g$ = 0.85，$g$ 95\% CI：0.02 至 1.69），FDR 校正后 $p$ = 0.126；训练自我效能组间差异极显著（+9.88 分，95\% CI：6.86 至 12.90，$p$ $<$ 0.001，$g$ = 2.70）。\textbf{用户接受度方面}，AI 组感知有用性（4.41 $\pm$ 0.20）、信任度（4.41 $\pm$ 0.20）与持续使用意愿（4.49 $\pm$ 0.27）均处于较高水平。质性访谈通过反思性主题分析归纳出 6 个核心主题与 5 类抗阻训练行为特征模式（积极接受型、安全依赖型、批判接受型、规则内化型、负荷导向型）。

\textbf{结论：}SportSci Pro 移动端计算机视觉系统在受控实验与真实训练场景中均展现出与专业线性位置传感器接近的测量精度，为低成本客观速度反馈方案下沉大众训练场景提供了测量学基础；数智化监控方案在依从受试者中具备扎实的执行可行性与安全性，但招募保留率与系统可用性未达预设推进标准，进入大规模确证性随机对照试验之前需针对性优化。4 周干预未在最大力量与纯向心爆发力上产生确证性组间优效，但在牵张缩短循环爆发力（CMJ）上捕获到值得进一步验证的方向性信号，在训练自我效能上产生了极显著的组间优势——该心理效应需在开放标签期望效应、模型残差非正态与量表天花板效应的三重审慎框架下解读。当前方案应定位为具有应用潜力、仍需大样本长周期验证的预试验性训练支持工具，向大众推广时不宜采用"一刀切"策略，而应根据用户训练经验实施分层自适应授权。

\textbf{关键词：}数智化监控；抗阻训练；基于速度的训练；预试验性研究；混合方法研究；随机对照预实验

% ========== 英文摘要 ==========
\clearpage
\begin{center}
\textbf{Abstract}
\end{center}
\addcontentsline{toc}{chapter}{Abstract}

\textbf{Objective:} Given that standardized repetitions-in-reserve (RIR)-based self-guidance has established scientific foundation and accessibility in resistance training, empirical evidence remains limited regarding whether a computer-vision-based mobile digital monitoring solution provides incremental adaptive benefit and implementation feasibility justifying its additional cost. This study employed a three-phase sequential design---instrument validity verification, pilot randomized controlled training comparison, and semi-structured interviews---to systematically evaluate the measurement foundation, intervention feasibility, exploratory training effects, and user acceptance of a self-developed mobile computer vision system (SportSci Pro) for digital monitoring of lower-body resistance training in young men, and to provide effect size estimates and process optimization parameters for a subsequent definitive randomized controlled trial (RCT).

\textbf{Methods:} In Study 1, 13 healthy young men with resistance training experience (independent of Study 2 sample) were recruited to synchronously collect mean concentric velocity (MCV) data during barbell back squats using SportSci Pro and a GymAware linear position transducer under controlled incremental-load laboratory conditions (60 paired observations); intraclass correlation coefficient (ICC), Bland-Altman analysis, and mean absolute error (MAE) were used to assess measurement agreement. Study 2 employed a two-arm parallel, open-label, stratified block-randomized controlled pilot design. Thirty-six healthy young men who completed randomization (AI-assisted group $n$ = 18, self-guidance group $n$ = 18) underwent a 4-week, twice-weekly lower-body resistance training intervention. During the squat main exercise, the AI-assisted group received velocity-loss-based automatic set termination, countermovement-jump (CMJ) readiness-based dynamic set reduction, and velocity-feedback load regulation, while the self-guidance group performed autonomous adjustments based on written training plans and RIR estimates. Primary feasibility outcomes included recruitment retention, training completion rate, subjective exertion, Hooper subjective recovery status, AI-assisted group velocity target hit rate, automatic set-reduction execution fidelity, System Usability Scale (SUS) score, and technology acceptance metrics. Exploratory outcomes included squat 1RM, CMJ, squat jump (SJ), and training self-efficacy. Concurrent SportSci Pro--GymAware paired data were collected in real training sessions for ecological validity supplement. Post-T1 testing, semi-structured interviews were conducted with 12 participants (AI group $n$ = 5, self-guidance group $n$ = 7); reflexive thematic analysis was used to elucidate subjective experiences and decision-making mechanisms.

\textbf{Results:} \textbf{For instrument validity}, Study 1 demonstrated excellent laboratory agreement between SportSci Pro and GymAware for MCV measurement (ICC = 0.940, 95\% CI: 0.854--0.975; MAE = 0.047 m/s; Bias = +0.002 m/s). In Study 2 real-world settings, rep-level paired analysis ($n$ = 43, 3 participants) yielded ICC = 0.991 (Bias = +0.0019 m/s); warm-up session-level paired analysis ($n$ = 88, all 11 AI-assisted group participants) yielded ICC = 0.987, but with a stable systematic positive bias of +0.0212 m/s (100\% of 88 pairs were positive). \textbf{For intervention feasibility}, of 36 randomized participants, 24 (AI group $n$ = 11, self-guidance group $n$ = 13) completed primary pre-post testing and entered per-protocol (PP) analysis. In the PP sample, training completion rate reached 98.4\% (189/192), high-adherence rate reached 100\% (24/24), AI group velocity target hit rate reached 80.2\% $\pm$ 6.9\%, automatic set-reduction execution fidelity reached 100\%, and no serious adverse events occurred. Seven of nine pre-specified progression criteria were met; however, the primary post-test completion rate (66.7\%, 24/36) and AI group SUS score (65.00 $\pm$ 4.61) did not meet the pre-specified thresholds of $\geq$80\% and $\geq$70 points, respectively. \textbf{For training execution and internal load}, the two groups were highly equivalent in total session load (AI-assisted 17651.8 $\pm$ 2198.6 kg vs. self-guidance 17461.5 $\pm$ 2792.1 kg, $p$ = 0.854), squat total load ($p$ = 0.988), overall mean session rating of perceived exertion (sRPE) (6.5 vs. 6.4), and individual Hooper means (20.4 vs. 19.8). The AI-assisted group completed significantly fewer total repetitions (100.3 vs. 105.5, $p$ = 0.035, Hedges' $g$ = $-$0.95), with a derived per-repetition mean mechanical load approximately 6.1\% higher. Hooper linear mixed-model analysis with Kenward-Roger correction revealed a non-significant group main effect ($p$ = 0.421), non-significant group-by-session interaction ($p$ = 0.945), and a significant session main effect ($p$ = 0.004); both groups exhibited a rise-then-fall temporal trajectory. \textbf{For exploratory adaptation outcomes}, the AI group showed higher point estimates than the self-guidance group for absolute 1RM (adjusted difference +2.43 kg, 95\% CI: $-$2.01 to 6.87; $g$ = 0.46), relative 1RM (+0.03 kg/kg; $g$ = 0.39), and SJ height (+0.39 cm; $g$ = 0.43), but all difference 95\% CIs crossed zero. For CMJ height, the adjusted difference was +2.16 cm (95\% CI: $-$0.01 to 4.33, $p$ = 0.051, $g$ = 0.85, $g$ 95\% CI: 0.02 to 1.69), with FDR-corrected $p$ = 0.126. Training self-efficacy showed a highly significant between-group difference (+9.88 points, 95\% CI: 6.86 to 12.90, $p$ $<$ 0.001, $g$ = 2.70). \textbf{For user acceptance}, AI group perceived usefulness (4.41 $\pm$ 0.20), trust (4.41 $\pm$ 0.20), and intention to continue use (4.49 $\pm$ 0.27) were all at relatively high levels. Reflexive thematic analysis of the interviews yielded 6 core themes and 5 resistance training behavioral pattern types (Active Acceptance, Safety Dependence, Critical Acceptance, Rule Internalization, and Load-Oriented).

\textbf{Conclusion:} The SportSci Pro mobile computer vision system demonstrated measurement precision approaching that of professional linear position transducer devices in both controlled laboratory and real training settings, providing a measurement foundation for the deployment of low-cost objective velocity feedback in general population resistance training. The digital monitoring protocol demonstrated solid implementation feasibility and safety among adherent participants; however, recruitment retention and system usability did not meet pre-specified progression criteria and require targeted optimization before advancing to a large-scale definitive RCT. The 4-week intervention did not produce confirmatory between-group superiority in maximal strength or pure concentric power, but captured a directionally consistent signal in stretch-shortening-cycle-based power (CMJ) warranting further verification, and produced a highly significant between-group advantage in training self-efficacy---this psychological effect must be interpreted within a three-fold caveat framework encompassing open-label expectation effects, non-normal model residuals, and potential scale ceiling effects. The current protocol should be positioned as a training support tool with application potential that requires large-sample, long-term validation. Rather than a one-size-fits-all rollout, a layered adaptive authorization strategy based on user training experience is recommended for general population deployment.

\addcontentsline{toc}{chapter}{Abstract}

\textbf{Key Words:} digital monitoring; resistance training; velocity-based training; pilot study; mixed methods research; randomized controlled pilot trial
